\section{What the checks establish and what they do not}
The operation family is part of the contract. A displacement result cannot
be inferred from a successful spacing edit, nor can a legacy positioned-block
test certify the newer normal-text movement path. Table~\ref{tab:policies}
summarizes the implemented distinctions.

\begin{table}[htbp]
\small
\begin{tabularx}{\linewidth}{@{}>{\raggedright\arraybackslash}p{.19\linewidth}YY@{}}
\toprule
Family & Acceptance evidence & Boundary\\
\midrule
Displacement / identity & Exact native source plan; target/page/glyph and non-text checks; complete same-renderer appearance & Bounded paint, target and page domain; not arbitrary document equivalence\\
Replacement / reflow & Exact native plan; visible target; non-target extracted-text multiplicities; non-text object count & Allows layout changes; does not prove unchanged unrelated geometry or graphic appearance\\
Parameter/style adapters & Exact source plan; eligible compile; independent readable inspection & Explicitly permits reflow; not the displacement preservation test\\
Legacy positioned blocks & Native compatibility plan and existing frontend acceptance & Separate compatibility path; not independently certified normal-text displacement\\
\bottomrule
\end{tabularx}
\caption{Distinct policies must remain distinct in experiments and claims.}
\label{tab:policies}
\end{table}

\subsection{Displacement and identity}
Current comparable target-geometry tolerance is 0.25 bp; unaffected-glyph
tolerance is 0.02 bp; non-text geometry uses $10^{-6}$ bp. A position-only move
must retain the relevant page structure and target correspondence. Identity
requires exact same-renderer pixels. The current fractional-displacement path
uses a baseline-derived render prediction rather than accepting approximate
candidate pixels or hiding a rectangular region around the target.

This strengthens checks against duplicated ink and unintended paint, but it
depends on correctly identifying the target and predicting its displacement.
Prediction and candidate inspection share PDFium. An error common to that
implementation can escape agreement. PDF.js supplies a different display/
selection path, but agreement between libraries is still not a formal proof
of document meaning, accessibility, metadata or every annotation.

The output contract is conditional on budgets and supported coordinates.
Rotated/cropped/scaled pages, annotations or complex paint can refuse. Long
documents have a bounded reading path; a reading/rendering pass is not evidence
that every page supports layout manipulation. A valid intended edit can lie
outside the supported subset and still be useful to a user, so its refusal is
reported as lost coverage rather than relabelled as an unsafe request.

\subsection{Reflow, parameters and legacy content}
Replacement may intentionally change word length, wrapping and page count.
The current reflow check preserves multiplicities of extracted non-target
text, checks non-text object count and requires visible target content. Those
predicates can miss a color, shape or positional change that preserves the
counts. The report does not extend the stronger displacement appearance
claim to this policy.

Rendered-reference agreement also cannot recover which equal source occurrence
was intended. The diagnostic in Section~\ref{sec:occurrence-reference} includes
deleting either of two adjacent identical paragraphs: the source patches differ,
while the entire rendered pages agree. Thus even a full appearance oracle needs
an independent source-target contract. This illustrates the already established
non-injective view/alignment problem; it is not a new general solution to it.

Semantic adapters likewise start with a bounded source grammar and exact
operation rederivation, then require eligible compilation/readable independent
inspection. Source-level image dimensions, table-column parameters and font/
spacing changes can cause legitimate global reflow. Their useful coverage and
failure cases need their own annotations. Legacy StudioBlock operations retain
a compatibility contract; the existence of independent verification elsewhere
does not imply that this older path has the same certification.

\subsection{Source preservation is not full TeX semantics}
Preserving all bytes outside a declared patch is valuable even when a change
would be invisible in a PDF. Comments, local line endings, definitions and
support files may matter to the author. Conversely, a tiny permitted source
change can have a large semantic effect through macro expansion or state.
Exact source preservation and rendered preservation therefore answer different
questions, and neither alone establishes the user's intended transformation.

The source-contract prototype checks a predeclared file set and baseline/
expected byte fingerprints. It can also bind a restricted fresh object-ID
slot in an otherwise exact UTF-8 template. It is not a TeX evaluator. A wrong
or incomplete contract is an oracle problem, not evidence that the application
is correct. The native-entry pilot uses manually declared existing-object
coordinate changes so its expected sources are not generated by the runtime
planner.

\subsection{Failure categories that must not collapse}
The evaluation distinguishes mapping ambiguity, unsupported grammar, source-plan
mismatch, compilation failure, independent-verifier refusal, resource limit,
stale/cancelled work and post-commit presentation failure. The location of a
refusal matters: a corrupt packet rejected by source rederivation supplies no
evidence about the final renderer's sensitivity to that packet.

Environmental failure is also different from a correct refusal. An engine
that cannot initialize or a missing font/package can prevent a meaningful
request from running. A clean developer-account build does not establish the
missing-dependency experience on another machine. Implemented setup guidance
helps expose these distinctions, but its native appearance and clean-account
use remain pending.

\subsection{Current evidence status}
Implemented mechanisms, executed tests and planned checks are separate.
Recorded engineering evidence includes automated source/protocol/controller
checks, real-engine corpora, older versioned Windows workflows and the component
pilots in Section~\ref{sec:development}. The current ordinary and diagnostic
application builds have not been opened or installed in this continuation,
at the owner's request. No final evaluation is frozen. No current evidence
establishes universal correctness, reference-machine latency, clean-machine
installation/uninstall or unassisted usability.
